\documentclass[12pt]{article}
\usepackage[margin=1in]{geometry}
\usepackage{hyperref}
\usepackage{graphicx}
\usepackage{listings}
\usepackage{xcolor}
\usepackage{booktabs}
\usepackage{cite}

% Code listing style
\lstset{
  basicstyle=\ttfamily\small,
  backgroundcolor=\color{gray!10},
  frame=single,
  breaklines=true,
  columns=fullflexible,
}

% ── Title block ────────────────────────────────────────────────────────────
\title{
  \textbf{Black-Box Fuzz Testing of the Graylayer Market Proxy API \\
  Using OpenAPI-Guided Test Generation}
}

\author{
  Asray Gopa \and
  Issmale Bekri \and
  Rohan Sah \and
  Akhil Pallem \and
  Peter Collins
}

\date{May 5, 2026}

% ───────────────────────────────────────────────────────────────────────────
\begin{document}

\maketitle

% ───────────────────────────────────────────────────────────────────────────
\section{Introduction}
\label{sec:intro}
\author{Asray Gopa}

Modern web services expose large attack surfaces through their APIs.
When developers write manual tests, they can only think of so many
inputs to try. This leaves a lot of edge cases untested, like weird
boundary values, missing fields, or malformed data. Fuzzing is a
testing technique that helps solve this problem by automatically
generating many different inputs and sending them to a program to
see if anything breaks.

In this project, we applied fuzzing to the Graylayer Market Proxy
API~\cite{graylayer_docs}, which is a service that pulls together
market data from Polymarket, Kalshi, Gemini, Coinbase, and
Forecastex. Since we do not have access to the source code, we treat
it as a black box. We first wrote an OpenAPI 3.0 specification based
on the published documentation, then used that spec to drive four
types of tests: schema-guided fuzzing with
Schemathesis~\cite{schemathesis_docs}, hand-crafted negative inputs,
differential tests that check cross-endpoint properties, and stateful
tests that verify list and detail endpoints stay consistent with each
other. The goal is to find crashes, bad error handling, and any
behavior that violates what the API is supposed to do.

% ───────────────────────────────────────────────────────────────────────────
\section{Previous and Existing Approaches}
\author{Peter Colins}
\label{sec:prev}

% TODO: fill in

% ───────────────────────────────────────────────────────────────────────────
\section{Technical Description}
\author{Issmale Bekri}
\label{sec:tech}

Schemathesis is a schema-guided API fuzzer built on top of the
Hypothesis property-based testing library~\cite{hypothesis_docs}.
It reads an OpenAPI specification and uses it to automatically
generate and send HTTP requests to a live API, then checks the
responses for crashes, invalid JSON, and slow responses. Since
we have no access to the Graylayer source code, we treat it as a
complete black box and rely on our hand-authored
\texttt{market\_proxy.yaml} to guide what inputs the fuzzer generates.

Before fuzzing begins, Schemathesis reads the
\texttt{market\_proxy.yaml} file and parses every endpoint, HTTP
method, parameter type, and response schema declared in it. This
gives the fuzzer a complete map of the API so it knows what valid
inputs look like for each endpoint before it starts generating test
cases.

After loading the spec, the fuzzing works as follows:

\begin{enumerate}
  \item Read the OpenAPI spec and collect all operations, which are
  each combination of an endpoint and an HTTP method.
  \item For each operation, use Hypothesis to generate many varied
  inputs based on the declared parameter types, constraints, and
  enums.
  \item Assemble each generated input into a complete HTTP request
  and send it to the live API with our API key in the
  \texttt{X-API-Key} header.
  \item Check the response against three assertions: the status
  code must not be a server crash (500 or 504), the body must be
  valid JSON if the Content-Type says so, and the response must
  not take longer than 10 seconds. Note that 502 and 503 responses
  are tolerated because the Graylayer docs document these as
  expected upstream errors from the proxy.
  \item If any assertion fails, Hypothesis shrinks the input down
  to the smallest version that still causes the failure.
  \item Move to the next operation and repeat.
\end{enumerate}

We can walk through a concrete example using the
\texttt{GET /api/v1/kalshi/markets} endpoint from our
\texttt{market\_proxy.yaml}. This endpoint declares a \texttt{limit}
parameter as an integer with \texttt{minimum: 1} and
\texttt{maximum: 1000}.

\begin{enumerate}
  \item Read the spec. The \texttt{limit} parameter is identified
  as an integer with minimum 1 and maximum 1000.

  \item Hypothesis generates test inputs. Say it generates
  \texttt{limit=100}, \texttt{limit=0}, \texttt{limit=-1}, and
  \texttt{limit=1001}. Schemathesis will normally generate many
  more, but we use these four as an example.

  \item Each input is assembled into a real HTTP GET request and
  sent to the live Graylayer endpoint with our API key in the
  \texttt{X-API-Key} header.

  \item The responses are checked. \texttt{limit=100} returns 200
  and passes all three assertions. \texttt{limit=0} and
  \texttt{limit=-1} are below the declared minimum so the server
  should return a 4xx rejection. \texttt{limit=1001} is above the
  declared maximum and should also be rejected. If any of these
  return a 500 error, the following assertion fails and the case
  is flagged:

\begin{lstlisting}[language=Python]
assert response.status_code < 500, (
    f"SERVER ERROR\n"
    f"  Endpoint : {case.method} {case.path}\n"
    f"  Status   : {response.status_code}\n"
    f"  Query    : {case.query}\n"
    f"  Body     : {response.text[:500]}"
)
\end{lstlisting}

  \item If a failure is found, Hypothesis shrinks the input. For
  example if \texttt{limit=-999} caused a 500 error, Hypothesis
  would try \texttt{limit=-1} and confirm that is the smallest
  input that still fails.

  \item Move to the next operation and repeat. Our spec covers
  around 40 endpoints across Coinbase, Polymarket, Kalshi, Gemini,
  and Forecastex.
\end{enumerate}

Beyond Schemathesis, we also run three additional test suites.
First, hand-crafted negative tests manually send malformed inputs
like null bytes, path traversal strings, and out-of-range values
to check how the API handles them. Second, differential tests check
that cross-endpoint properties hold, for example that the best bid
is always less than or equal to the best ask. Third, stateful tests
verify that every market ID returned by a list endpoint can also be
fetched from the matching detail endpoint. Together these four
techniques give much broader coverage than any single approach on
its own.

% ───────────────────────────────────────────────────────────────────────────
\section{Evaluation}
\author{Rohan Sah}
\label{sec:eval}

% TODO: fill in
To evaluate the technique described in Section~\ref{sec:tech}, we ran
Schemathesis against two production Graylayer services: the Market Proxy
API and the companion Orderbook History API. We authored a separate
\texttt{openapi.yaml} for each service from the public documentation,
giving the fuzzer a combined map of 55 unique operations spanning the
Coinbase, Polymarket, Kalshi, Gemini, Bitmex, and Forecastex venues.
Tests were executed from the Schemathesis 3.x command-line tool with
the default \texttt{--checks all} preset, an \texttt{X-API-Key} header
attached to every request, and a per-operation cap of 100 generated
examples. Each run took roughly twelve minutes on a single thread, and
all responses were captured to a JSON Lines log so that individual
events could be replayed and inspected after the fact.

Across both services the fuzzer recorded 3{,}064 events. Schemathesis
classifies an ``event'' as a single generated request together with the
outcome of the three assertions described in Section~\ref{sec:tech}:
the status code must be below 500, the body must parse as JSON when
the \texttt{Content-Type} header advertises it, and the body must
conform to the response schema declared in the spec. Of the 3{,}064
events, 2{,}403 were flagged as unexpected by at least one of the
three checks, while the remaining cases passed cleanly. The
unexpected-event rate was uneven across venues: Coinbase and Gemini
endpoints failed checks on roughly one in ten requests, whereas Kalshi
and Bitmex endpoints failed checks on more than half, indicating that
the published documentation for those two venues drifts noticeably
from the live response shape.

The most serious findings were 32 HTTP 500 responses and 5
high-severity \texttt{negative:crash} events. The crashes were not
distributed evenly. All five clustered in three endpoints:
\texttt{GET /api/v1/kalshi/snapshots},
\texttt{GET /api/v1/kalshi/deltas}, and
\texttt{GET /api/v1/bitmex/tickers}. After Hypothesis shrunk the
failing inputs, the smallest reproducer for the Kalshi snapshots
crash was a \texttt{from\_ts} value larger than \texttt{to\_ts}; the
service returned a 500 with an empty body rather than a 4xx
validation error. The Bitmex tickers crash was triggered by a
\texttt{symbol} query parameter containing an unescaped percent sign,
which appears to fail an internal URL-decode step inside the proxy
layer before any venue-specific validation runs. These are exactly
the kind of edge cases a manually written test suite is unlikely to
include, and they were surfaced with no instrumentation of the
target binary.

The remaining unexpected events fell into two large buckets. First,
57 events were rate-limit signals (HTTP 429) that the spec did not
declare as a possible response, so Schemathesis correctly flagged
them as schema contract violations even though the underlying
behavior is desirable. Second, 28 events were transport-layer
failures---read timeouts and TCP resets---most of which clustered
around the higher-volume Polymarket history endpoints under
sustained load. The bulk of the remaining 2{,}300+ unexpected events
were genuine schema drift: nullable fields declared as required,
timestamp fields returned as integers where the spec promised
strings, and array fields returned as \texttt{null} on empty
results.

The technique's main advantages were visible immediately: with no
access to the Graylayer source and only a hand-authored spec, we
generated several thousand request/response pairs and isolated five
reproducible crashes in under thirty minutes of wall time. Hypothesis's
shrinking gave us a minimal failing input for every defect, which
makes filing a report straightforward. The principal disadvantage is
that coverage is bounded by the spec: any endpoint, parameter, or
response field omitted from \texttt{openapi.yaml} is invisible to
the fuzzer, and stateful flows that require a sequence of calls
(for example, creating a session token before reading user-scoped
data) are out of reach without additional hooks. Schema-guided
fuzzing is therefore best understood as a low-effort first pass for
any third-party API where source access is unavailable: it surfaces
crashes and contract drift quickly, but should be paired with manual
review and authenticated end-to-end tests before a system is
declared safe.

% ───────────────────────────────────────────────────────────────────────────
\section{Summary and Recommendation}
\author{Akhil Pallem}
\label{sec:summary}

In this project, we evaluated schema-guided black-box API fuzzing applied to two
Graylayer services. The Market Proxy API and the Orderbook History API using
Schemathesis alongside hand crafted negative, differential, and stateful test suites.
The fuzzer exercised 55 unique endpoints and recorded 3,064 events, checking each response against the three assertions described in Section~\ref{sec:tech}: status code
below 500, valid JSON body, and schema contract compliance. Despite these checks,
32 HTTP 500 responses slipped through, 2,403 events were flagged as unexpected, and
5 high severity \texttt{negative:crash} findings were concentrated in the Kalshi
snapshots, deltas, and Bitmex tickers endpoints demonstrating that even a
well-documented API can harbor fault surfaces that manual testing would not reach.

Schema-guided fuzzing proved effective at surfacing server crashes, rate-limit
violations (57 events), and transport errors (28 events) with minimal manual effort
beyond authoring the OpenAPI specifications introduced in Section~\ref{sec:intro}.
Coverage is inherently bounded by the accuracy of those specs, and behaviors not
described in the YAML remain invisible to the fuzzer---a core limitation of the
black-box approach. We recommend this technique for security assessments of
third-party APIs where source access is unavailable, and as a low-effort first pass
before manual penetration testing. To build a more resilient testing lifecycle, it should
be paired with static schema linting, authenticated end-to-end tests, and targeted
manual review of the crash cases identified here.

% TODO: fill in

% ───────────────────────────────────────────────────────────────────────────
\bibliographystyle{IEEEtran}
\begin{thebibliography}{9}

\bibitem{graylayer_docs}
Graylayer, ``Graylayer API Documentation,''
\url{https://docs.graylayer.tech}, accessed Apr.\ 2026.

\bibitem{schemathesis_docs}
Schemathesis Contributors, ``Schemathesis: Property-Based Testing for API Schemas,''
\url{https://schemathesis.readthedocs.io}, accessed Apr.\ 2026.

\bibitem{hypothesis_docs}
D. R. MacIver \textit{et al.}, ``Hypothesis: A new approach to property-based testing,''
\textit{Journal of Open Source Software}, vol.\ 4, no.\ 43, p.\ 1891, 2019.

\bibitem{fuzzing_usenix}
A. Gutmann, ``Fuzzing: The state of the art,''
\textit{USENIX ;login:}, vol.\ 41, no.\ 2, 2016.
[Online]. Available: \url{https://www.usenix.org/system/files/login/articles/login_summer16_03_gutmann.pdf}

\bibitem{afl}
M. Zalewski, ``American Fuzzy Lop (AFL),''
\url{https://github.com/google/AFL}, accessed Apr.\ 2026.

\end{thebibliography}

\end{document}
